% Options for packages loaded elsewhere
\PassOptionsToPackage{unicode}{hyperref}
\PassOptionsToPackage{hyphens}{url}
\PassOptionsToPackage{dvipsnames,svgnames,x11names}{xcolor}
%
\documentclass[
  11pt,
]{article}
\usepackage{amsmath,amssymb}
\usepackage{iftex}
\ifPDFTeX
  \usepackage[T1]{fontenc}
  \usepackage[utf8]{inputenc}
  \usepackage{textcomp} % provide euro and other symbols
\else % if luatex or xetex
  \usepackage{unicode-math} % this also loads fontspec
  \defaultfontfeatures{Scale=MatchLowercase}
  \defaultfontfeatures[\rmfamily]{Ligatures=TeX,Scale=1}
\fi
\usepackage{lmodern}
\ifPDFTeX\else
  % xetex/luatex font selection
\fi
% Use upquote if available, for straight quotes in verbatim environments
\IfFileExists{upquote.sty}{\usepackage{upquote}}{}
\IfFileExists{microtype.sty}{% use microtype if available
  \usepackage[]{microtype}
  \UseMicrotypeSet[protrusion]{basicmath} % disable protrusion for tt fonts
}{}
\makeatletter
\@ifundefined{KOMAClassName}{% if non-KOMA class
  \IfFileExists{parskip.sty}{%
    \usepackage{parskip}
  }{% else
    \setlength{\parindent}{0pt}
    \setlength{\parskip}{6pt plus 2pt minus 1pt}}
}{% if KOMA class
  \KOMAoptions{parskip=half}}
\makeatother
\usepackage{xcolor}
\usepackage[margin=1in]{geometry}
\usepackage{color}
\usepackage{fancyvrb}
\newcommand{\VerbBar}{|}
\newcommand{\VERB}{\Verb[commandchars=\\\{\}]}
\DefineVerbatimEnvironment{Highlighting}{Verbatim}{commandchars=\\\{\}}
% Add ',fontsize=\small' for more characters per line
\newenvironment{Shaded}{}{}
\newcommand{\AlertTok}[1]{\textcolor[rgb]{1.00,0.00,0.00}{\textbf{#1}}}
\newcommand{\AnnotationTok}[1]{\textcolor[rgb]{0.38,0.63,0.69}{\textbf{\textit{#1}}}}
\newcommand{\AttributeTok}[1]{\textcolor[rgb]{0.49,0.56,0.16}{#1}}
\newcommand{\BaseNTok}[1]{\textcolor[rgb]{0.25,0.63,0.44}{#1}}
\newcommand{\BuiltInTok}[1]{\textcolor[rgb]{0.00,0.50,0.00}{#1}}
\newcommand{\CharTok}[1]{\textcolor[rgb]{0.25,0.44,0.63}{#1}}
\newcommand{\CommentTok}[1]{\textcolor[rgb]{0.38,0.63,0.69}{\textit{#1}}}
\newcommand{\CommentVarTok}[1]{\textcolor[rgb]{0.38,0.63,0.69}{\textbf{\textit{#1}}}}
\newcommand{\ConstantTok}[1]{\textcolor[rgb]{0.53,0.00,0.00}{#1}}
\newcommand{\ControlFlowTok}[1]{\textcolor[rgb]{0.00,0.44,0.13}{\textbf{#1}}}
\newcommand{\DataTypeTok}[1]{\textcolor[rgb]{0.56,0.13,0.00}{#1}}
\newcommand{\DecValTok}[1]{\textcolor[rgb]{0.25,0.63,0.44}{#1}}
\newcommand{\DocumentationTok}[1]{\textcolor[rgb]{0.73,0.13,0.13}{\textit{#1}}}
\newcommand{\ErrorTok}[1]{\textcolor[rgb]{1.00,0.00,0.00}{\textbf{#1}}}
\newcommand{\ExtensionTok}[1]{#1}
\newcommand{\FloatTok}[1]{\textcolor[rgb]{0.25,0.63,0.44}{#1}}
\newcommand{\FunctionTok}[1]{\textcolor[rgb]{0.02,0.16,0.49}{#1}}
\newcommand{\ImportTok}[1]{\textcolor[rgb]{0.00,0.50,0.00}{\textbf{#1}}}
\newcommand{\InformationTok}[1]{\textcolor[rgb]{0.38,0.63,0.69}{\textbf{\textit{#1}}}}
\newcommand{\KeywordTok}[1]{\textcolor[rgb]{0.00,0.44,0.13}{\textbf{#1}}}
\newcommand{\NormalTok}[1]{#1}
\newcommand{\OperatorTok}[1]{\textcolor[rgb]{0.40,0.40,0.40}{#1}}
\newcommand{\OtherTok}[1]{\textcolor[rgb]{0.00,0.44,0.13}{#1}}
\newcommand{\PreprocessorTok}[1]{\textcolor[rgb]{0.74,0.48,0.00}{#1}}
\newcommand{\RegionMarkerTok}[1]{#1}
\newcommand{\SpecialCharTok}[1]{\textcolor[rgb]{0.25,0.44,0.63}{#1}}
\newcommand{\SpecialStringTok}[1]{\textcolor[rgb]{0.73,0.40,0.53}{#1}}
\newcommand{\StringTok}[1]{\textcolor[rgb]{0.25,0.44,0.63}{#1}}
\newcommand{\VariableTok}[1]{\textcolor[rgb]{0.10,0.09,0.49}{#1}}
\newcommand{\VerbatimStringTok}[1]{\textcolor[rgb]{0.25,0.44,0.63}{#1}}
\newcommand{\WarningTok}[1]{\textcolor[rgb]{0.38,0.63,0.69}{\textbf{\textit{#1}}}}
\usepackage{longtable,booktabs,array}
\usepackage{calc} % for calculating minipage widths
% Correct order of tables after \paragraph or \subparagraph
\usepackage{etoolbox}
\makeatletter
\patchcmd\longtable{\par}{\if@noskipsec\mbox{}\fi\par}{}{}
\makeatother
% Allow footnotes in longtable head/foot
\IfFileExists{footnotehyper.sty}{\usepackage{footnotehyper}}{\usepackage{footnote}}
\makesavenoteenv{longtable}
\setlength{\emergencystretch}{3em} % prevent overfull lines
\providecommand{\tightlist}{%
  \setlength{\itemsep}{0pt}\setlength{\parskip}{0pt}}
\setcounter{secnumdepth}{-\maxdimen} % remove section numbering
\ifLuaTeX
  \usepackage{selnolig}  % disable illegal ligatures
\fi
\IfFileExists{bookmark.sty}{\usepackage{bookmark}}{\usepackage{hyperref}}
\IfFileExists{xurl.sty}{\usepackage{xurl}}{} % add URL line breaks if available
\urlstyle{same}
\hypersetup{
  pdfauthor={Stefano Alimonti --- stefalimonti@gmail.com},
  colorlinks=true,
  linkcolor={blue},
  filecolor={Maroon},
  citecolor={Blue},
  urlcolor={blue},
  pdfcreator={LaTeX via pandoc}}

\author{Stefano Alimonti --- stefalimonti@gmail.com}
\date{March 2026}

\begin{document}

{
\hypersetup{linkcolor=black}
\setcounter{tocdepth}{2}
\tableofcontents
}
\section{The Carry--Dirichlet Bridge: Stopping-Time Series and
L²(s,χ₄)}\label{the-carrydirichlet-bridge-stopping-time-series-and-luxb2sux3c7ux2084}

\textbf{Author:} Stefano Alimonti \textbf{Affiliation:} Independent Researcher
\textbf{Date:} March 2026

\begin{center}\rule{0.5\linewidth}{0.5pt}\end{center}

\newpage

\subsection{Abstract}\label{abstract}

We construct a Dirichlet series from the stopping-time decomposition of the
binary carry chain for D-odd multiplication and study its relation to Dirichlet
L-functions. The construction proceeds through an exact algebraic bridge: Witt
p-adic correction vectors biject onto classical carry vectors (Proposition 1),
the carry and Witt transition operators are conjugate via an explicit
permutation intertwiner (Proposition 2), and a χ₄-weighted character channel
extracts the resolvent identity to machine precision (Proposition 3).

The resulting carry-Dirichlet channel μ\_χ(K,s) = Σ\_τ u\_mix(τ) χ(n(τ))
n(τ)\^{}\{−s\} converges uniformly on compact subsets of Re(s) \textgreater{} 1
as K → ∞, with geometric Cauchy contraction governed by the Diaconis--Fulman
spectral gap 1/2 (Theorem 1). For the character χ₄ (mod 4), the limit satisfies

\begin{quote}
μ\_χ₄\^{}(∞)(s) = c̃ · (1 + 3\^{}\{−s\}) · L(s, χ₄)² + ε(s),
\end{quote}

with holdout relative error \textbar ε\textbar/\textbar μ\textbar{} ≈ 2.5\% and
cross-K stability below 10⁻⁵ (Theorem 2, numerical). The correction factor (1 +
3\^{}\{−s\}) removes exactly one power of the Euler factor at p = 3, and the L²
structure reflects the Dirichlet convolution arising from the product x · y. The
correction factor F(s) = μ/(cL) is meromorphic and zero-free in the tested strip
region {[}0.3, 0.7{]} × {[}0, 15{]} (Proposition 5). A low-complexity corrector
law extends to χ₅ and χ₈ with modified exponent, while χ₃ (conductor 3) is an
outlier.

These results establish a quantitative bridge between the Markov dynamics of
binary carries and the multiplicative structure of Dirichlet L-functions. They
do not imply zero-transfer theorems or results toward the Riemann Hypothesis.

\begin{center}\rule{0.5\linewidth}{0.5pt}\end{center}

\newpage

\subsection{1. Introduction}\label{introduction}

\subsubsection{1.1 Motivation}\label{motivation}

Papers {[}P1{]} and {[}E{]} provide strong numerical evidence that the cascade
sector ratio R(K) of D-odd binary multiplication converges to −π = −4L(1, χ₄)
(Conjecture 1 of {[}P1{]}; 4.0-digit exponential Richardson extrapolation
through K = 21), connecting the carry chain to the Dirichlet character χ₄. The
stopping-time decomposition of {[}P1, §8.2a{]} expresses R(K) as a sum over
depths τ, with each term factoring as P(τ \textbar{} sector) · E{[}val
\textbar{} τ, sector{]}.

A natural question arises: does this stopping-time structure encode deeper
information about L-functions beyond the single value L(1, χ₄)? In this paper we
form the Dirichlet series

\[\mu_\chi(K,s) = \sum_{\tau \geq \tau_0} u_{\text{mix}}(\tau;\, K)\, \chi(n(\tau))\, n(\tau)^{-s},\]

where n(τ) = 2(τ − τ₀) + 1 maps stopping depths to odd integers and u\_mix = ω ·
u₁₀ − u₀₀ combines the sector contributions weighted by the count ratio ω =
N₁₀/N₀₀. We show that this series converges to a limit on Re(s) \textgreater{} 1
and that the limit is accurately modeled by a corrected square of L(s, χ).

\subsubsection{1.2 Main results}\label{main-results}

\textbf{Algebraic bridge (§2).} We construct a five-step exact bridge from Witt
p-adic arithmetic to the resolvent decomposition of R(K):

\begin{enumerate}
\def\labelenumi{\arabic{enumi}.}
\tightlist
\item
  Carry vectors biject onto Witt correction vectors (Proposition 1).
\item
  Carry and Witt transition operators are conjugate via an explicit permutation
  (Proposition 2).
\item
  A χ₄-like character channel produces identical coefficients on both sides
  (Proposition 3).
\item
  The construction extends to the nonstationary D-odd multiplication setting
  (§2.4).
\item
  The resolvent identity Σ\_τ u(τ) = μ\_ab holds to machine precision
  (Proposition 4).
\end{enumerate}

\textbf{Convergence (§3).} The family \{μ\_χ(K, s)\}\_K converges uniformly on
compact subsets of Re(s) \textgreater{} 1 as K → ∞, with

\[\sup_{s \in \mathcal{K}} |\mu_\chi(K, s) - \mu_\chi(K', s)| \leq C \rho^{\min(K,K')}\]

for tested compacts K ⊂ \{Re(s) \textgreater{} 1\}, with ρ ≈ 0.5--0.6 consistent
with the Diaconis--Fulman spectral gap (Theorem 1).

\textbf{The L² formula (§4).} For χ₄, the high-K limit satisfies

\[\mu_{\chi_4}^{(\infty)}(s) \approx \tilde{c} \cdot (1 + 3^{-s}) \cdot L(s, \chi_4)^2\]

on Re(s) \textgreater{} 1, with holdout relative error ≈ 2.5\% (Theorem 2,
numerical). Two equivalent formulations of this identity are:

\begin{itemize}
\tightlist
\item
  Euler-tail form: μ ≈ c · L(s, χ₄) · ∏\_\{p\textgreater3\} (1 − χ₄(p)
  p\textsuperscript{\{−s\})}\{−1\}
\item
  Corrected-L² form: μ ≈ c̃ · (1 + 3\^{}\{−s\}) · L(s, χ₄)²
\end{itemize}

These coincide on Re(s) \textgreater{} 1 because (1 + 3\^{}\{−s\}) · L² = L ·
∏\_\{p≥5\} (1 − χ₄(p) p\textsuperscript{\{−s\})}\{−1\}. The Euler-tail
truncation to p ≤ 401 introduces a discrepancy below 3 × 10⁻³ on tested grids.

\textbf{Zero-free structure (§5).} The correction factor F(s) = μ(s)/(c · L(s,
χ₄)) is meromorphic in the tested box {[}0.3, 0.7{]} × {[}0, 15{]} with poles at
the zeros of L(s, χ₄) and no zeros of its own (Proposition 5). This implies that
μ\^{}(∞) is analytic and zero-free in that region.

\textbf{Cross-character extension (§6).} A low-complexity corrector law

\[\mu_\chi(s) \approx c \cdot (1 - \chi(p_0) p_0^{-s})^k \cdot L(s, \chi)^2\]

extends to χ₅ (k = 2, gain +46\%), χ₈ (k = 2, gain +71\%), with p₀ = 3 in all
cases. The character χ₃ is an outlier: since χ₃(3) = 0 (the conductor divides
p₀), the corrector is inert.

\subsubsection{1.3 Scope and non-claims}\label{scope-and-non-claims}

This paper is a contribution to experimental mathematics. The L² formula
(Theorem 2) is a numerical observation with controlled error bounds, not a
proof. We do not claim:

\begin{itemize}
\tightlist
\item
  An exact identity μ = c̃ · (1 + 3\^{}\{−s\}) · L² (the 2.5\% residual is
  genuine).
\item
  A zero-transfer theorem from the carry chain to L-function zeros.
\item
  Any result toward the Riemann Hypothesis.
\end{itemize}

\subsubsection{1.4 Notation}\label{notation}

Throughout this paper: K denotes the bit-length of the multiplicands, D = 2K − 1
the product length, b = 2 the base, χ\_q a primitive real Dirichlet character
modulo q, and L(s, χ) = Σ\_\{n≥1\} χ(n) n\^{}\{−s\} the Dirichlet L-function.

\begin{center}\rule{0.5\linewidth}{0.5pt}\end{center}

\newpage

\subsection{2. The Algebraic Bridge}\label{the-algebraic-bridge}

This section constructs the exact operator-level bridge from Witt p-adic
arithmetic to the stopping-time resolvent of the carry chain.

\subsubsection{2.1 Witt vectors and carry
corrections}\label{witt-vectors-and-carry-corrections}

For integers x, y ∈ {[}0, 2\^{}K), the classical carry vector c = (c₀, c₁,
\ldots, c\_D) of x + y satisfies c₀ = 0 and c\_\{j+1\} = ⌊(x\_j + y\_j +
c\_j)/2⌋. The Witt correction vector δ of the same addition is defined by

\[\delta_j = c_j - 2 c_{j+1},\]

which takes values in \{−2, −1, 0, +1\}.

\textbf{Proposition 1} (Carry--Witt bijection). \emph{The map c ↦ δ is a
bijection between classical carry vectors and Witt correction vectors for
truncated base-2 addition. Moreover, the canonical truncated Witt embedding is a
ring homomorphism for addition and multiplication.}

Verification: exact for all 2\^{}\{2K\} pairs at K ≤ 10 (experiment L01).

\subsubsection{2.2 Operator intertwiner}\label{operator-intertwiner}

Define the augmented carry state z\_j = (c\_j, c\_\{j+1\}) ∈ \{(0,0), (0,1),
(1,0), (1,1)\} and the Witt correction state δ\_j = c\_j − 2c\_\{j+1\} ∈ \{0,
+1, −1, −2\}. The map φ: z\_j ↦ δ\_j is bijective on the four-element state
space.

Let T\_z and T\_d denote the empirical transition matrices on z-states and
δ-states respectively, and let P be the 4 × 4 permutation matrix induced by φ.

\textbf{Proposition 2} (Operator conjugacy). \emph{T\_d = P T\_z P\^{}\{−1\}.
The carry-side and Witt-side dynamics are exactly conjugate.}

Verification: exact conjugacy at K = 7, \ldots, 12 (experiment L02).

\subsubsection{2.3 Character channel}\label{character-channel}

Define the χ₄-like weight vector on δ-states as w\_d = (0, 0, +1, −1),
corresponding to the even/odd decomposition of the correction values. The
pullback to z-states is w\_z = w\_d · P.

The character channel coefficient at depth τ is

\[\text{ch}(\tau) = w_z^\top \cdot v_z(\tau),\]

where v\_z(τ) is the empirical z-state distribution at depth τ.

\textbf{Proposition 3} (Channel identity). \emph{For all tested K and depths τ:
the carry-side and Witt-side character channel coefficients coincide, ch\_z(τ) =
ch\_d(τ). After removing the stationary component, the transient decays at the
half-rate 1/2.}

Verification: exact identity at K = 7, \ldots, 12 (experiment L03).

\subsubsection{2.4 Nonstationary D-odd
extension}\label{nonstationary-d-odd-extension}

For D-odd binary multiplication (both factors K-bit, product exactly D = 2K − 1
bits), the transition operators T(τ) are depth-dependent. We construct per-depth
operators from large random samples of D-odd pairs and verify:

\begin{itemize}
\tightlist
\item
  The intertwiner P conjugates T\_z(τ) and T\_d(τ) at every depth.
\item
  The character channel equality holds per-depth.
\end{itemize}

In the sampled tests, the per-depth spectral gap \textbar λ₂(τ)\textbar{} ranges
from 0.31 to 0.63 across K = 7, \ldots, 12, with a decreasing trend in K. At
small K (K = 7, 8), the gap can exceed the Diaconis--Fulman reference rate
0.500; by K ≥ 10, all observed gaps fall below 0.500 (experiment L04, sampled).

\subsubsection{2.5 Resolvent decomposition}\label{resolvent-decomposition}

The stopping-time decomposition expresses the cascade mean as

\[\mu_{ab} = \sum_{\tau} u_{ab}(\tau), \qquad u_{ab}(\tau) = P(\tau \mid ab) \cdot E[\text{val} \mid \tau, ab].\]

\textbf{Proposition 4} (Resolvent identity). \emph{For all K = 7, \ldots, 12 and
both sectors ab ∈ \{00, 10\}: the stopping-time sum reproduces the direct
cascade mean to machine precision (error ≤ 6 × 10⁻¹⁷).}

The R(K) values from exact enumeration at K = 7, \ldots, 12 match the
independent E45 high-K profiles to full available precision (experiment L05).

\subsubsection{2.6 Spectral gap}\label{spectral-gap}

\textbf{Proposition 6} (Diaconis--Fulman bound, numerical). \emph{The numerical
evidence strongly supports that conditioning on D-odd does not change the
dominant spectral gap: the two-step survival product q(τ)q(τ+1) converges to 1/4
= (1/2)², consistent with the Diaconis--Fulman rate for the conditioned chain. A
fully rigorous proof requires showing that the Doob h-transform of the carry
chain (conditioning on D-odd) has spectral radius ≤ 1/2 for the restricted
chain.} (Experiment L06.)

\begin{center}\rule{0.5\linewidth}{0.5pt}\end{center}

\newpage

\subsection{3. Uniform Convergence}\label{uniform-convergence}

\subsubsection{3.1 The carry-Dirichlet
channel}\label{the-carry-dirichlet-channel}

For a primitive real Dirichlet character χ modulo q, define the carry-Dirichlet
channel:

\[\mu_\chi(K, s) = \sum_{\tau \geq \tau_0} u_{\text{mix}}(\tau;\, K)\, \chi(n(\tau))\, n(\tau)^{-s},\]

where τ₀ = 3, n(τ) = 2(τ − τ₀) + 1, and u\_mix(τ; K) = ω(K) · u₁₀(τ; K) − u₀₀(τ;
K).

The coefficients u\_mix(τ) decay geometrically with rate ≈ 1/2 (the
Diaconis--Fulman rate), so the series converges absolutely on all of ℂ for each
fixed K.

\subsubsection{3.2 Convergence diagnostics}\label{convergence-diagnostics}

We evaluate μ\_χ(K, s) on a compact grid in Re(s) \textgreater{} 1 for K ∈ \{7,
\ldots, 11\} (exact enumeration) and K ∈ \{19, 20, 21, 999\} (E45 profiles and
Richardson extrapolation). The proxy K = 999 is obtained by Richardson
extrapolation of the K = 20 and K = 21 profiles.

\textbf{Theorem 1} (Uniform Cauchy convergence, numerical). \emph{For each
tested character χ ∈ \{χ₃, χ₄, χ₅, χ₈\}:}

\emph{(a) The high-K contraction ratios satisfy}

\[\frac{\delta_{20}}{\delta_{19}} \leq 0.61, \qquad \frac{\delta_{21}}{\delta_{20}} \leq 0.61,\]

\emph{where δ\_K = sup\_s \textbar μ\_χ(K, s) − μ\_χ(999, s)\textbar.}

\emph{(b) The Cauchy chain is monotone: ‖μ₁₉ − μ₂₀‖ ≥ ‖μ₂₀ − μ₂₁‖ ≥ ‖μ₂₁ −
μ₉₉₉‖.}

\emph{(c) An empirical geometric envelope δ\_K ≤ C ρ\^{}K holds with ρ
\textless{} 0.95, with zero violations on tested K values.}

These properties hold uniformly across all four tested characters (experiment
L07).

\begin{center}\rule{0.5\linewidth}{0.5pt}\end{center}

\newpage

\subsection{4. The L² Formula}\label{the-luxb2-formula}

\subsubsection{4.1 Proportionality to L}\label{proportionality-to-l}

On Re(s) \textgreater{} 1, the ratio μ\_χ₄(K, s) / L(s, χ₄) converges to an
approximately constant value c(K) as K grows. At K = 999, the maximum relative
deviation of the ratio from its mean is ≈ 8\% on holdout grids (experiment L08).

\subsubsection{4.2 Euler structure of the
correction}\label{euler-structure-of-the-correction}

The correction factor F(s) = μ(s)/(c · L(s, χ₄)), normalized at a reference
point s₀ = 1.70, is well-approximated by a partial Euler product over odd
primes:

\[F^N(s) \approx \prod_{5 \leq p \leq 401} \frac{1}{1 - \chi_4(p)\, p^{-s}} \bigg/ \text{(value at } s_0\text{)},\]

with hold-MRE = 4.3\%, a 54\% improvement over the constant-residual baseline,
at zero free parameters. Cross-K stability is extraordinary: rel-std = 8 × 10⁻⁶
across K ∈ \{19, 20, 21, 999\} (experiment L08).

\textbf{Methodological caveat.} The candidate (Pc, Pmax) was selected by
minimizing MRE on the same hold-grid reported above, so the L08 hold-MRE figure
carries a selection-bias component. The independent bootstrap audit in L09 and
the cross-K stability metric provide the true unbiased validation.

\subsubsection{4.3 The corrected-L² form}\label{the-corrected-luxb2-form}

\textbf{Theorem 2} (Corrected L² law, numerical). \emph{For χ₄, the high-K
carry-Dirichlet channel on Re(s) \textgreater{} 1 satisfies}

\[\mu_{\chi_4}^{(\infty)}(s) = \tilde{c} \cdot (1 + 3^{-s}) \cdot L(s, \chi_4)^2 + \varepsilon(s),\]

\emph{with:}

\begin{itemize}
\tightlist
\item
  \emph{Holdout relative error: max \textbar ε\textbar/\textbar μ\textbar{} ≈
  2.5\% on a 15-point grid disjoint from the fitting grid.}
\item
  \emph{Cross-K stability: relative standard deviation of hold-MRE across K ∈
  \{19, 20, 21, 999\} is below 10⁻⁵.}
\end{itemize}

The factor (1 + 3\^{}\{−s\}) = (1 − χ₄(3) · 3\^{}\{−s\}) is the inverse of the
Euler factor at p = 3. Its presence indicates that the Dirichlet convolution
arising from the product x · y doubles all Euler factors except the one at p =
3. The algebraic equivalence

\[(1 + 3^{-s}) \cdot L(s, \chi_4)^2 = L(s, \chi_4) \cdot \prod_{p \geq 5} (1 - \chi_4(p)\, p^{-s})^{-1}\]

explains why the Euler-tail and corrected-L² formulations produce near-identical
fits (hold-MRE 2.57\% vs 2.50\%) on Re(s) \textgreater{} 1 (experiment L09).

\subsubsection{4.4 Interpretation: Dirichlet
convolution}\label{interpretation-dirichlet-convolution}

The L² structure has a natural arithmetic explanation. The carry-Dirichlet
channel arises from the product x · y of two K-bit integers. Each factor
contributes independently to the divisibility structure of the product, so the
natural Dirichlet series weight is the convolution χ ∗ χ, whose generating
function is L(s, χ)². The correction at p = 3 reflects a residual interaction
between the binary carry dynamics and the smallest odd prime: since 3 = 2 + 1,
the carry overflow at p = 3 is partially absorbed by the binary structure,
preventing a full doubling of the Euler factor.

\subsubsection{4.5 Residual audit}\label{residual-audit}

A comprehensive audit (experiment L09) confirms:

\begin{itemize}
\tightlist
\item
  The corrected-L² model outperforms all tested alternatives (L alone, L²,
  uncorrected L · E\_tail).
\item
  Bootstrap resampling (1200 iterations) confirms the advantage of (1 +
  3\^{}\{−s\}) · L² over the Euler-tail formulation is small but consistent
  (median advantage 6 × 10⁻⁴, 95\% CI entirely negative).
\item
  Leave-one-prime-out ablation identifies p = 5 as the dominant contributor:
  removing it degrades hold-MRE by +0.115.
\item
  Sign-shuffled controls yield p\_emp = 0.045 (below the conventional 5\%
  threshold, confirming non-random structure).
\end{itemize}

\begin{center}\rule{0.5\linewidth}{0.5pt}\end{center}

\newpage

\subsection{5. Zero-Free Structure}\label{zero-free-structure}

\subsubsection{5.1 Zeros of μ in the critical
strip}\label{zeros-of-ux3bc-in-the-critical-strip}

For K ∈ \{19, 20, 21, 999\}, the function μ\_χ₄(K, s) has no zeros in the box
{[}0.3, 0.7{]} × {[}0, 15{]} of the critical strip, as determined by the
argument principle (winding number computation on rectangular contours).

For comparison, L(s, χ₄) has three zeros in the same box, at approximately s ≈
0.5 + 6.02i, 0.5 + 10.24i, 0.5 + 12.99i.

\subsubsection{5.2 Zero-free correction
factor}\label{zero-free-correction-factor}

\textbf{Proposition 5} (F zero-free). \emph{The correction factor F(s) = μ(s)/(c
· L(s, χ₄)) has winding index −3 on the boundary of {[}0.3, 0.7{]} × {[}0,
15{]}, decomposing as Z − P = 0 − 3 where Z = 0 (no zeros) and P = 3 (three
poles, inherited from the three zeros of L(s, χ₄)). This is consistent with F
being meromorphic and zero-free in the tested strip region.}

Within the tested box, this implies that μ\^{}(∞) is itself analytic and
zero-free: the carry chain does not develop zeros at the L-function zero
positions in this region. Extension to the full critical strip would require
either larger contour tests or an analytic argument; neither is available at
present.

\subsubsection{5.3 Critical-line mirroring}\label{critical-line-mirroring}

On the critical line Re(s) = 1/2, the function \textbar μ\_χ₄(999, 1/2 +
it)\textbar{} exhibits local minima near the imaginary parts of the L-function
zeros (t ≈ 5.91, 10.31, 12.81 vs the exact values 6.02, 10.24, 12.99). These
minima reflect the harmonic content of the carry coefficients but do not
converge to zeros as K increases: the depth-of-minima ratios S\_j(K) remain
approximately flat across K = 19, \ldots, 999.

This confirms the structural limitation: the carry chain encodes the amplitude
of L(s, χ₄) on Re(s) \textgreater{} 1 (via the L² formula) but cannot reproduce
the phase cancellations that create zeros in the strip.

\begin{center}\rule{0.5\linewidth}{0.5pt}\end{center}

\newpage

\subsection{6. Cross-Character Analysis}\label{cross-character-analysis}

\subsubsection{6.1 Local-corrector extension}\label{local-corrector-extension}

The L² formula generalizes to other primitive real characters via a
low-complexity corrector:

\[\mu_\chi(s) \approx c \cdot (1 - \chi(p_0)\, p_0^{-s})^k \cdot L(s, \chi)^2,\]

where p₀ is the smallest odd prime with χ(p₀) ≠ 0, and k is a small integer. A
discrete scan over (p, k) with p ∈ \{3, 5, \ldots, 29\}, k ∈ \{−2, −1, 1, 2\}
yields:

\begin{longtable}[]{@{}llllll@{}}
\toprule\noalign{}
Character & q & Best p₀ & Best k & Hold-MRE & Gain vs L² \\
\midrule\noalign{}
\endhead
\bottomrule\noalign{}
\endlastfoot
χ₄ & 4 & 3 & 1 & 0.025 & +84\% \\
χ₈ & 8 & 3 & 2 & 0.075 & +71\% \\
χ₅ & 5 & 3 & 2 & 0.247 & +46\% \\
χ₃ & 3 & --- & --- & (no gain) & −6\% \\
\end{longtable}

For χ₄, the best model coincides with the predicted corrector (p₀ = 3, k = 1);
no scan is needed. Cross-K stability is below 10⁻⁵ for all improving characters
(experiment L10).

\subsubsection{6.2 The χ₃ outlier}\label{the-ux3c7ux2083-outlier}

The character χ₃ (conductor 3) fails to improve because χ₃(3) = 0: the corrector
at p₀ = 3 is trivially 1. No single-prime correction at larger primes (p = 5, 7,
\ldots, 29) compensates. This is a structural limitation: when the conductor
divides the natural corrector prime, the mechanism is inert.

\subsubsection{6.3 Mechanism analysis}\label{mechanism-analysis}

A/B controls (experiment L11, verdict: MIXED/INCONCLUSIVE) provide partial
evidence on the mechanism underlying the corrector pattern (p₀ = 3 with
character-dependent k):

\begin{itemize}
\tightlist
\item
  \textbf{Not driven by a single factor}: removing the n\^{}\{−s\} weight
  preserves the k-pattern, but removing the character channel or scrambling the
  τ → n map degrades it. However, shuffled-χ and scrambled-n controls do not
  fully separate the contributions (some gates fail).
\item
  \textbf{Not reducible to raw p-adic frequencies}: the distribution of v₃(xy)
  among D-odd pairs does not isolate p = 3 as dominant. The correction appears
  to be an analytic-structural phenomenon involving the interaction of the
  character channel, the stopping-time map, and the Dirichlet weight, though a
  complete mechanistic derivation remains open.
\end{itemize}

\begin{center}\rule{0.5\linewidth}{0.5pt}\end{center}

\newpage

\subsection{7. Discussion}\label{discussion}

\subsubsection{7.1 Summary}\label{summary}

The carry-Dirichlet bridge establishes a quantitative connection between three
mathematical objects:

\begin{enumerate}
\def\labelenumi{\arabic{enumi}.}
\tightlist
\item
  The \textbf{Markov dynamics} of binary carry propagation (spectral gap 1/2,
  Diaconis--Fulman).
\item
  The \textbf{stopping-time decomposition} of D-odd multiplication (resolvent
  identity, cascade values).
\item
  The \textbf{Dirichlet L-function} L(s, χ₄) and its square L²(s, χ₄).
\end{enumerate}

The connection takes the form μ\^{}(∞)(s) ≈ c̃ · (1 + 3\^{}\{−s\}) · L²(s, χ₄) on
Re(s) \textgreater{} 1, with ≈ 2.5\% residual. The L² structure arises from the
Dirichlet convolution inherent in multiplication, and the (1 + 3\^{}\{−s\})
correction reflects the partial absorption of the p = 3 Euler factor by the
binary carry dynamics.

\subsubsection{7.2 Relation to companion
papers}\label{relation-to-companion-papers}

\begin{itemize}
\tightlist
\item
  \textbf{{[}P1{]}} conjectures R(∞) = −π = −4L(1, χ₄) (Conjecture 1, supported
  by 4.0-digit numerical evidence) and develops the stopping-time decomposition.
  The present paper extends the analysis from a single value (s = 1) to a
  function of s, revealing the L² structure.
\item
  \textbf{{[}E{]}} proves R = −π conditionally on LMH via the shifted resolvent.
  The L² formula provides independent confirmation and a deeper structural
  explanation.
\item
  \textbf{{[}H{]}} shows that carry corrections do not contribute information
  about zeta zeros beyond the Euler product. The zero-free property of §5 is
  consistent with this finding.
\item
  \textbf{{[}Frobenius{]}} constructs the Witt-carry bridge at odd primes p = 3,
  7. Propositions 1--2 of the present paper establish the analogous bridge at p
  = 2.
\end{itemize}

\subsubsection{7.3 Open problems}\label{open-problems}

\begin{enumerate}
\def\labelenumi{\arabic{enumi}.}
\item
  \textbf{Exact identity or approximation?} Is there a closed form for the 2.5\%
  residual? The correction F(s) may involve higher-order Euler factors or
  finitely many additional primes.
\item
  \textbf{Analytic computation of c̃.} The proportionality constant c̃ should be
  derivable from the carry chain structure. Numerically c̃ ≈ 0.025 (for the (1 +
  3\^{}\{−s\}) · L² formulation).
\item
  \textbf{Why p = 3?} The correction involves p = 3 = 2 + 1 exclusively. A
  precise connection between the base b = 2 and the corrector prime p = b + 1
  may hold in generality.
\item
  \textbf{Multi-prime extension.} Carry chains at base p = 3, 5, 7, \ldots{}
  would provide the local factors at each prime. The product of all local
  factors should converge to L(s, χ) itself (single power, not squared),
  eliminating the need for the Euler-tail correction.
\item
  \textbf{Conductor dependence of k.} The exponent k = 1 for q = 4 and k = 2 for
  q = 5, 8 may follow from the structure of (ℤ/qℤ)* and its interaction with the
  base b = 2. A structural derivation would replace the current discrete scan.
\item
  \textbf{Universal character law.} The failure of χ₃ (conductor 3) suggests
  that a universal law requires treating the case gcd(q, p₀) \textgreater{} 0
  separately. Whether a two-prime corrector (p₀ = 3, p₁ = 5) rescues χ₃ is
  untested.
\end{enumerate}

\begin{center}\rule{0.5\linewidth}{0.5pt}\end{center}

\newpage

\subsection{8. Computational Appendix}\label{computational-appendix}

\subsubsection{8.1 Experiment index}\label{experiment-index}

\begin{longtable}[]{@{}
  >{\raggedright\arraybackslash}p{(\columnwidth - 4\tabcolsep) * \real{0.2667}}
  >{\raggedright\arraybackslash}p{(\columnwidth - 4\tabcolsep) * \real{0.4333}}
  >{\raggedright\arraybackslash}p{(\columnwidth - 4\tabcolsep) * \real{0.3000}}@{}}
\toprule\noalign{}
\begin{minipage}[b]{\linewidth}\raggedright
Script
\end{minipage} & \begin{minipage}[b]{\linewidth}\raggedright
Description
\end{minipage} & \begin{minipage}[b]{\linewidth}\raggedright
Section
\end{minipage} \\
\midrule\noalign{}
\endhead
\bottomrule\noalign{}
\endlastfoot
\texttt{\_shared.py} & Shared utilities: E45 data loading, profiles,
L-functions, characters & --- \\
\texttt{data/} & Bundled high-K profile data (E45, E162) & --- \\
\texttt{L01\_witt\_carry\_identity.py} & Witt--carry bijection and ring
homomorphism & §2.1 \\
\texttt{L02\_operator\_intertwiner.py} & Operator conjugacy via permutation
intertwiner & §2.2 \\
\texttt{L03\_character\_channel.py} & χ₄ character channel and transient rate &
§2.3 \\
\texttt{L04\_dodd\_nonstationary.py} & D-odd extension with per-depth operators
& §2.4 \\
\texttt{L05\_analytic\_resolvent.py} & Resolvent identity and R(K) validation &
§2.5 \\
\texttt{L06\_diaconis\_fulman\_bound.py} & Spectral gap confirmation under D-odd
conditioning & §2.6 \\
\texttt{L07\_uniform\_convergence.py} & Cauchy convergence diagnostics & §3 \\
\texttt{L08\_euler\_tail\_identity.py} & Euler-tail structure of the correction
& §4.2 \\
\texttt{L09\_residual\_audit.py} & Bootstrap, ablation, controls; L² equivalence
& §4.3--§4.5 \\
\texttt{L10\_local\_corrector.py} & Cross-character corrector scan & §6.1 \\
\texttt{L11\_mechanism\_controls.py} & A/B controls for mechanism analysis &
§6.3 \\
\end{longtable}

\subsubsection{8.2 Data sources}\label{data-sources}

\begin{itemize}
\tightlist
\item
  \textbf{Exact enumeration} (K = 7, \ldots, 12): full D-odd pair enumeration
  via \texttt{L05\_analytic\_resolvent.py}.
\item
  \textbf{High-K profiles} (K = 19, 20, 21): per-position cascade profiles from
  E45 ({[}E{]}, experiments E45\_high\_K\_profiles.c). These are the same data
  used in {[}P1{]} and {[}E{]}.
\item
  \textbf{K = 999 proxy}: Richardson extrapolation of K = 20 and K = 21
  profiles, following {[}P1, §9.2{]}.
\end{itemize}

\subsubsection{8.3 Reproduction}\label{reproduction}

High-K profile data (E45) and R(K) history (E162) are in
\texttt{experiments/data/}. Shared utilities (profile building, L-function
evaluation, character definitions) are in \texttt{experiments/\_shared.py}.

\begin{Shaded}
\begin{Highlighting}[]
\ExtensionTok{pip}\NormalTok{ install numpy mpmath}
\BuiltInTok{cd}\NormalTok{ experiments/}
\ExtensionTok{python}\NormalTok{ L01\_witt\_carry\_identity.py}
\ExtensionTok{python}\NormalTok{ L02\_operator\_intertwiner.py}
\ExtensionTok{python}\NormalTok{ L03\_character\_channel.py}
\ExtensionTok{python}\NormalTok{ L04\_dodd\_nonstationary.py}
\ExtensionTok{python}\NormalTok{ L05\_analytic\_resolvent.py }\AttributeTok{{-}{-}k{-}min}\NormalTok{ 7 }\AttributeTok{{-}{-}k{-}max}\NormalTok{ 12}
\ExtensionTok{python}\NormalTok{ L06\_diaconis\_fulman\_bound.py}
\ExtensionTok{python}\NormalTok{ L07\_uniform\_convergence.py}
\ExtensionTok{python}\NormalTok{ L08\_euler\_tail\_identity.py}
\ExtensionTok{python}\NormalTok{ L09\_residual\_audit.py}
\ExtensionTok{python}\NormalTok{ L10\_local\_corrector.py}
\ExtensionTok{python}\NormalTok{ L11\_mechanism\_controls.py}
\end{Highlighting}
\end{Shaded}

\subsubsection{8.4 Gate summary}\label{gate-summary}

All quantitative claims in this paper are supported by explicit pass/fail gates
with pre-specified thresholds. The complete gate definitions are documented in
the experiment scripts.

\begin{longtable}[]{@{}
  >{\raggedright\arraybackslash}p{(\columnwidth - 4\tabcolsep) * \real{0.3333}}
  >{\raggedright\arraybackslash}p{(\columnwidth - 4\tabcolsep) * \real{0.2857}}
  >{\raggedright\arraybackslash}p{(\columnwidth - 4\tabcolsep) * \real{0.3810}}@{}}
\toprule\noalign{}
\begin{minipage}[b]{\linewidth}\raggedright
Claim
\end{minipage} & \begin{minipage}[b]{\linewidth}\raggedright
Gate
\end{minipage} & \begin{minipage}[b]{\linewidth}\raggedright
Result
\end{minipage} \\
\midrule\noalign{}
\endhead
\bottomrule\noalign{}
\endlastfoot
Witt--carry bijection (Prop 1) & exact match all pairs K ≤ 10 & PASS \\
Operator conjugacy (Prop 2) & ‖T\_d − P T\_z P⁻¹‖ \textless{} 10⁻¹⁴ & PASS \\
Channel identity (Prop 3) & \(\lvert ch_z - ch_d \rvert < 10^{-14}\) & PASS \\
Resolvent identity (Prop 4) & \(\lvert \mu_{res} - \mu_{dir} \rvert < 10^{-16}\)
& PASS \\
Cauchy convergence (Thm 1) & contraction ratios ≤ 0.75, envelope violation = 0 &
PASS \\
L² formula (Thm 2) & hold-MRE ≤ 0.03, cross-K rel-std ≤ 10⁻² & PASS \\
F zero-free in tested box (Prop 5) & winding index Z = 0, P = N\_L & PASS \\
Cross-character (§6) & ≥ 3 characters with ≥ 40\% gain, rel-std ≤ 0.02 & PASS \\
\end{longtable}

\begin{center}\rule{0.5\linewidth}{0.5pt}\end{center}

\newpage

\subsection{References}\label{references}

\begin{itemize}
\tightlist
\item
  {[}P1{]} S. Alimonti, \emph{Pi from Pure Arithmetic: A Spectral Phase
  Transition in the Binary Carry Bridge}, 2026.
\item
  {[}E{]} S. Alimonti, \emph{The Trace Anomaly of Binary Multiplication}, 2026.
\item
  {[}A{]} S. Alimonti, \emph{Spectral Theory of Carries}, 2026.
\item
  {[}B{]} S. Alimonti, \emph{Carry Polynomials and the Euler Product}, 2026.
\item
  {[}H{]} S. Alimonti, \emph{Carry Polynomials and the Partial Euler Product
  (Control)}, 2026.
\item
  {[}Frobenius{]} S. Alimonti, \emph{Frobenius Eigenvalues and Gauss Sums from
  Witt Carries}, 2026.
\item
  P. Diaconis, J. Fulman, \emph{Carries, shuffling, and symmetric functions},
  Advances in Applied Mathematics 43 (2009), 176--196.
\end{itemize}

\end{document}
